\documentclass[11pt]{ctexart}
\usepackage[a4paper,margin=1in]{geometry}
\usepackage{graphicx}
\usepackage{xcolor}
\usepackage{hyperref}
\definecolor{refgray}{gray}{0.45}
\title{\textbf{市场最新观点汇总}\\[0.4em]{\large 通胀粘性、政策转向与结构性需求切换主导市场，AI与能源基建重构产业链估值逻辑。}}
\author{KC桌面 自动生成}
\date{260526}
\begin{document}
\maketitle
\section*{一页摘要}
\begin{itemize}
  \item 美联储可能从降息转向加息，市场对政策方向逆转的定价不足。
  \item 中国4月经济数据偏弱但非全面塌陷，地产出现企稳迹象，人民币国际化与半导体自主化构成结构性变量。
  \item AI投资对GDP的拉动被统计口径低估，电力链路重构催生模拟半导体长期增长。
  \item 全球储能与化工供应链正经历永久性收缩，中国技术生态优势深化。
  \item 代币化RWA逆势增长，传统金融机构主导的资产上链正在改变金融基础设施。
\end{itemize}
\section{宏观与利率}
\textbf{美联储政策框架面临重新校准，市场对加息风险的定价不足；中国宏观数据偏弱但政策信号转向稳定。}

\begin{itemize}
  \item 美联储此前175个基点的降息可能被证明是过度保险，当前通胀加速迫使政策制定者考虑加息选项，最鸽派的官员也开始提及加息可能性。
  \item 密歇根大学5月长期通胀预期上修至3.9\%，已回到2022年恐慌水平，消费者不相信当前政策能拉回通胀至2\%目标。
  \item 中国4月工业增加值、社零、固投全面低于预期，但部分弱来自全球能源冲击和财政支出节奏放缓，并非需求端全面塌陷。
  \item 中国地产高频数据出现企稳迹象，一线城市新房和二手房价格环比上涨，但价格预期指数连续走弱，市场可能进入量稳价跌阶段。
  \item 法国国债面临夏季供给高峰、政治日程与评级行动三重逆风，45\%的净供给集中在6-8月，穆迪负面展望可能触发评级行动。
\end{itemize}
\begin{figure}[htbp]
\centering
\includegraphics[width=0.88\linewidth]{figures/fig_001.jpg}
\caption{Figure 1 - Figure 1: EM USD Aggregate 10s30s}
\end{figure}
\begin{figure}[htbp]
\centering
\includegraphics[width=0.88\linewidth]{figures/fig_002.jpg}
\caption{Figure 2 - Figure 2: UST 10s30s yield curve slope}
\end{figure}
\begin{figure}[htbp]
\centering
\includegraphics[width=0.88\linewidth]{figures/fig_148.jpg}
\caption{Exhibit 1 - Exhibit 1: Primary GFA sold last week was \$-4\textbackslash{}\%\$ wow and \$-5\textbackslash{}\%\$ yoy in c.75 cities}
\end{figure}
\begin{figure}[htbp]
\centering
\includegraphics[width=0.88\linewidth]{figures/fig_149.jpg}
\caption{Exhibit 2 - Exhibit 2: Primary GFA sold YTD on average was -14\% yoy in c.75 cities, and -14\%/-48\% vs. 2024/2023 level}
\end{figure}
\begin{figure}[htbp]
\centering
\includegraphics[width=0.88\linewidth]{figures/fig_150.jpg}
\caption{Exhibit 3 - Exhibit 3: Secondary GFA sold last week was +1\% wow and +11\% yoy in c.20 cities Average weekly volume of secondary property sales}
\end{figure}
\begin{figure}[htbp]
\centering
\includegraphics[width=0.88\linewidth]{figures/fig_151.jpg}
\caption{Exhibit 4 - Exhibit 4: Secondary GFA sold YTD was -1\% yoy in c.20 cities, while +23\%/+6\% vs. 2024/ secondary volume sold vs. 2022-25}
\end{figure}
\begin{figure}[htbp]
\centering
\includegraphics[width=0.88\linewidth]{figures/fig_152.jpg}
\caption{Exhibit 5 - Exhibit 5: Average CSI was -1.6pp wow and +2.3pp yoy Weekly Centraline Salesman Index (CSI) tracker in 5 cities}
\end{figure}
\begin{figure}[htbp]
\centering
\includegraphics[width=0.88\linewidth]{figures/fig_153.jpg}
\caption{Exhibit 6 - Exhibit 6: Average CAI was -0.4pp wow and -4.1pp yoy Weekly Centraline Seller Asking Index (CAI) tracker in 6 cities}
\end{figure}
{\small\color{refgray}\textbf{References:} [R001] 市场真正低估的不是通胀，而是美联储可能被迫“反向保险”加息; [R002] 市场真正低估的不是需求，而是政策信号的再定价; [R010] 法国债券的夏季风暴尚未被市场定价; [R012] 市场真正低估的不是成交量，而是价格预期的自我强化}

\section{AI与算力}
\textbf{AI投资对GDP的拉动被统计口径严重低估，电力链路重构成为比芯片更关键的瓶颈。}

\begin{itemize}
  \item AI相关支出2026年将拉动真实资本开支增长约3.3个百分点，但GDP统计因将半导体视为中间投入品，仅计入0.1个百分点。
  \item 单机架功耗从传统15kW跃升至下一代1.5MW以上，现有48V/12V供电架构无法支撑，行业正转向800V直流架构。
  \item 到2030年，仅AI数据中心从机架到芯片的模拟半导体市场将从79亿美元扩张至270亿美元，价值向最靠近加速器的环节集中。
  \item 华为提出Tau缩放定律，从几何缩放转向时间缩放，通过LogicFolding实现Cell-to-Cell堆叠，提供可预测的十年性能提升路线。
  \item 人形机器人在仓库场景的经济账已算得过来，当前投资回报期约3年，规模量产后可缩短至1年以内。
\end{itemize}
\begin{figure}[htbp]
\centering
\includegraphics[width=0.88\linewidth]{figures/fig_036.jpg}
\caption{Exhibit 1 - Exhibit 1: After Stripping Out Increases in Equipment Imports and Other Special Factors, We Estimate That “Underlying” Domestic Capex Growth Picked Up to 5\% in 2026Q1 from 2.4\% in 2025 QoQ Annualized Real Nonresidential Fixed Inv}
\end{figure}
\begin{figure}[htbp]
\centering
\includegraphics[width=0.88\linewidth]{figures/fig_037.jpg}
\caption{Exhibit 2 - \# AI Investment We expect AI-related spending to continue boosting equipment and structures investment in 2026 as companies press ahead with the infrastructure buildout. Over time, the impact on software and R\&D should also become more visible as enterprise ad}
\end{figure}
\begin{figure}[htbp]
\centering
\includegraphics[width=0.88\linewidth]{figures/fig_038.jpg}
\caption{Exhibit 4 - Exhibit 3: We Estimate That AI-Related Spending Will Boost True Capex Growth by About 3.3pp and Measured Capex Growth by 2.8pp in 2026, With the Gap Reflecting Undermeasurement of Semiconductor Investment AI Spending Contribution}
\end{figure}
\begin{figure}[htbp]
\centering
\includegraphics[width=0.88\linewidth]{figures/fig_039.jpg}
\caption{Exhibit 4 - Exhibit 4: We Estimate That AI-Related Spending Will Boost True GDP Growth by 0.3pp but Measured GDP Growth by Only 0.1pp in 2026, With the Gap Reflecting Undermeasurement of Semiconductor Investment and AI-Related Intellectual Pro}
\end{figure}
\begin{figure}[htbp]
\centering
\includegraphics[width=0.88\linewidth]{figures/fig_115.jpg}
\caption{exhibit 1 - Multiple bottlenecks are emerging in the multi-year AI infrastructure buildout – memory, optics, leading-edge logic wafers, advanced substrates, etc. – but among the greatest constraints to scaling is arguably power. \# Higher compute density is translating to }
\end{figure}
\begin{figure}[htbp]
\centering
\includegraphics[width=0.88\linewidth]{figures/fig_116.jpg}
\caption{Exhibit 5 - Exhibit 5: Cumulative data center power demand globally is expected to increase from \textbackslash{}\textasciitilde{}100GW today to almost 300GW by CY30 Data center power demand by region (GW)}
\end{figure}
\begin{figure}[htbp]
\centering
\includegraphics[width=0.88\linewidth]{figures/fig_117.jpg}
\caption{Exhibit 6 - Exhibit 6: Layering on top of conventional server power demand, AI is catalyzing electricity consumption to new heights Data center electricity consumption (TWh)}
\end{figure}
\begin{figure}[htbp]
\centering
\includegraphics[width=0.88\linewidth]{figures/fig_118.jpg}
\caption{Exhibit 7 - Exhibit 7: We estimate AI accelerator demand implies that GWs installed per year for compute could 4x from 15GW in CY25 to 60GW in CY30, a cumulative 233GWs deployed during the period Implied GWs deployed across major accelerator p}
\end{figure}
\begin{figure}[htbp]
\centering
\includegraphics[width=0.88\linewidth]{figures/fig_119.jpg}
\caption{Exhibit 8 - Exhibit 8: Today, power from the grid undergoes multiple conversions from high voltage AC to medium voltage AC and then multiple low voltage DCs before reaching the accelerator Traditional 48V/54V power distribution architecture}
\end{figure}
\begin{figure}[htbp]
\centering
\includegraphics[width=0.88\linewidth]{figures/fig_120.jpg}
\caption{Exhibit 9 - Exhibit 9: 800V DC architecture minimizes energy losses and supports higher reliability by making fewer AC to DC and DC to DC conversions along with reducing system component complexity Existing power flow architecture vs. next gen}
\end{figure}
\begin{figure}[htbp]
\centering
\includegraphics[width=0.88\linewidth]{figures/fig_125.jpg}
\caption{Exhibit 15 - Exhibit 15: We expect the AI TAM for analog semis to grow 28\% 5-year CAGR through CY30 to \textbackslash{}\$27bn from \textbackslash{}\$7.9bn in CY25. Average content per rack could grow from mid-\textbackslash{}\$20K to over \textbackslash{}\$60K by CY30 Analog semi data center and strategic po}
\end{figure}
\begin{figure}[htbp]
\centering
\includegraphics[width=0.88\linewidth]{figures/fig_126.jpg}
\caption{Exhibit 17 - Exhibit 17: While analog ICs will continue to be the largest market, we believe SiC and GaN will enjoy the most share gains AI analog semi TAM by device type}
\end{figure}
\begin{figure}[htbp]
\centering
\includegraphics[width=0.88\linewidth]{figures/fig_127.jpg}
\caption{Exhibit 18 - Exhibit 18: TXN enjoys dominant share in the AI TAM but we flag that discrete suppliers ON and Infineon gain significant share Analog semi vendor AI market share}
\end{figure}
\begin{figure}[htbp]
\centering
\includegraphics[width=0.88\linewidth]{figures/fig_128.jpg}
\caption{Exhibit 20 - Exhibit 20: The data center analog TAM at \textbackslash{}\$7.6bn CY25 growing to \textbackslash{}\$25bn by CY30, a \textbackslash{}\textasciitilde{}28\% CAGR Analog semi AI data center TAM}
\end{figure}
\begin{figure}[htbp]
\centering
\includegraphics[width=0.88\linewidth]{figures/fig_049.jpg}
\caption{Exhibit 7 - EXHIBIT 1: Humanoid robots work in warehouses: material handling and sorting Industrial robot arm operating in a factory setting with yellow structural components (no visible text or symbols) Agility Robotics}
\end{figure}
\begin{figure}[htbp]
\centering
\includegraphics[width=0.88\linewidth]{figures/fig_050.jpg}
\caption{EXHIBIT 2 - EXHIBIT 3: Illustration of logistic process ```mermaid graph TD}
\end{figure}
\begin{figure}[htbp]
\centering
\includegraphics[width=0.88\linewidth]{figures/fig_051.jpg}
\caption{EXHIBIT 4 - EXHIBIT 4: Workflow in warehouses ```mermaid graph TD}
\end{figure}
\begin{figure}[htbp]
\centering
\includegraphics[width=0.88\linewidth]{figures/fig_052.jpg}
\caption{EXHIBIT 5 - EXHIBIT 5: Warehouses with different automation levels Overhead view of a large warehouse with workers sorting packages on pallets and shipping containers (no visible text or signage) Low automation level}
\end{figure}
\begin{figure}[htbp]
\centering
\includegraphics[width=0.88\linewidth]{figures/fig_053.jpg}
\caption{EXHIBIT 7 - EXHIBIT 7: Global warehouse automation solution market size and penetration rate}
\end{figure}
\begin{figure}[htbp]
\centering
\includegraphics[width=0.88\linewidth]{figures/fig_054.jpg}
\caption{EXHIBIT 8 - EXHIBIT 8: Warehouse automation market breakdown by solutions in 2024 Global warehouse automation market in 2024 (total market size USD65bn)}
\end{figure}
\begin{figure}[htbp]
\centering
\includegraphics[width=0.88\linewidth]{figures/fig_055.jpg}
\caption{EXHIBIT 9 - EXHIBIT 9: Barriers to adopting and implementing warehouse automation What do you consider to be the biggest barriers to adopting and implementing warehouse AS/RS system in your organization?}
\end{figure}
\begin{figure}[htbp]
\centering
\includegraphics[width=0.88\linewidth]{figures/fig_056.jpg}
\caption{EXHIBIT 10 - EXHIBIT 10: US warehousing \& storage labor demand has grown over the past decade}
\end{figure}
\begin{figure}[htbp]
\centering
\includegraphics[width=0.88\linewidth]{figures/fig_057.jpg}
\caption{EXHIBIT 11 - EXHIBIT 11: The U.S. warehousing and storage workforce is shifting toward fulfillment-driven roles}
\end{figure}
\begin{figure}[htbp]
\centering
\includegraphics[width=0.88\linewidth]{figures/fig_058.jpg}
\caption{EXHIBIT 12 - EXHIBIT 12: The U.S.: Rising wages have increased labor costs in warehouses}
\end{figure}
\begin{figure}[htbp]
\centering
\includegraphics[width=0.88\linewidth]{figures/fig_059.jpg}
\caption{EXHIBIT 13 - EXHIBIT 13: China logistics and warehouse employees China logistics \& warehouse employees breakdown in 2023}
\end{figure}
\begin{figure}[htbp]
\centering
\includegraphics[width=0.88\linewidth]{figures/fig_060.jpg}
\caption{EXHIBIT 14 - EXHIBIT 14: China logistics and warehouse employees' annual wages Average annual wage of China logistics \& warehouse employees (urban non-private units only) in 2024}
\end{figure}
\begin{figure}[htbp]
\centering
\includegraphics[width=0.88\linewidth]{figures/fig_061.jpg}
\caption{EXHIBIT 15 - EXHIBIT 15: Payback analysis of humanoid robots in the U.S. warehouses EXHIBIT 16: The U.S.: Price declines in humanoid robots will shorten the payback period United States: Payback period of humanoid robots in warehouses (Assume 75}
\end{figure}
\begin{figure}[htbp]
\centering
\includegraphics[width=0.88\linewidth]{figures/fig_062.jpg}
\caption{EXHIBIT 19 - EXHIBIT 19: China: Price declines in humanoid robots will shorten the payback period China: Payback period of humanoid robots in warehouses (Assume 75\% human worker efficiency)}
\end{figure}
\begin{figure}[htbp]
\centering
\includegraphics[width=0.88\linewidth]{figures/fig_067.jpg}
\caption{EXHIBIT 1 - EXHIBIT 1: Huawei's tau scaling roadmap τ-Scaling Roadmap: Sustainable PPDC Evolution}
\end{figure}
{\small\color{refgray}\textbf{References:} [R004] 市场真正低估的不是AI支出规模，而是它如何被GDP统计“吃掉”; [R006] 人形机器人在仓库的回报期已缩短至3年，规模量产将使其低于1年; [R007] 市场低估了“Tau定律”对半导体国产化的重塑力; [R011] 市场真正低估的不是AI芯片，而是从电网到芯片的电力链路重构}

\section{能源与大宗}
\textbf{化工与储能供应链正经历永久性收缩，中国技术生态优势深化，地缘政治风险消退后的油价仍具粘性。}

\begin{itemize}
  \item 即便中东冲突结束，上游化工品供给曲线已发生不可逆抬升，长期聚烯烃价格预测上调5\%-15\%，部分老旧装置在经济上不值得修复。
  \item 全球储能竞争从产能竞赛进入技术生态与供应链主权阶段，福特工厂将韩国设备全部替换为CATL技术，转向LFP体系。
  \item 霍尔木兹海峡若重启，油价仍将粘性停留在100美元/桶以上，油轮短缺和补库需求将延续瓶颈。
  \item 阳光电源拿下阿联酋7.5GWh储能订单，采用8小时充电/16小时放电模式，中国企业在高温环境下的技术优势凸显。
  \item 固态电池加速推进，赣锋锂业开始小批量生产500Wh/kg锂金属固态电池，蜂巢能源计划提前量产半固态电池。
\end{itemize}
\begin{figure}[htbp]
\centering
\includegraphics[width=0.88\linewidth]{figures/fig_168.jpg}
\caption{Exhibit 4 - EXHIBIT 2: In the US, olefin prices have rallied since the onset of the Iran conflict and show no signs of easing. Olefin US-Gulf prices index - Rebased to 100}
\end{figure}
\begin{figure}[htbp]
\centering
\includegraphics[width=0.88\linewidth]{figures/fig_169.jpg}
\caption{EXHIBIT 2 - EXHIBIT 2: In the US, olefin prices have rallied since the onset of the Iran conflict and show no signs of easing. Olefin US-Gulf prices index - Rebased to 100}
\end{figure}
\begin{figure}[htbp]
\centering
\includegraphics[width=0.88\linewidth]{figures/fig_170.jpg}
\caption{EXHIBIT 3 - EXHIBIT 3: Chinese olefin prices have surged since the onset of the Iran conflict and began stabilizing in April.}
\end{figure}
\begin{figure}[htbp]
\centering
\includegraphics[width=0.88\linewidth]{figures/fig_171.jpg}
\caption{EXHIBIT 6 - EXHIBIT 6: LYB also have made the case for lasting changes to the petrochemical industry in their 1Q26 presentation \# Middle East war creating structural shifts in economics}
\end{figure}
\begin{figure}[htbp]
\centering
\includegraphics[width=0.88\linewidth]{figures/fig_172.jpg}
\caption{EXHIBIT 6 - EXHIBIT 6: LYB also have made the case for lasting changes to the petrochemical industry in their 1Q26 presentation \# Middle East war creating structural shifts in economics LYB is well-positioned to benefit from strengthened cost}
\end{figure}
\begin{figure}[htbp]
\centering
\includegraphics[width=0.88\linewidth]{figures/fig_173.jpg}
\caption{Exhibit 7 - EXHIBIT 7: Key commodity chemicals are typically used as intermediates for other chemical processes: (1: Ethylene) World 2023 Ethylene Demand}
\end{figure}
\begin{figure}[htbp]
\centering
\includegraphics[width=0.88\linewidth]{figures/fig_174.jpg}
\caption{EXHIBIT 8 - EXHIBIT 8: Key commodity chemicals are typically used as intermediates for other chemical processes: (2: Propylene) World 2023 Propylene Demand}
\end{figure}
\begin{figure}[htbp]
\centering
\includegraphics[width=0.88\linewidth]{figures/fig_175.jpg}
\caption{EXHIBIT 9 - EXHIBIT 9: Polyethylene uses vary by grade (part 1: high density) Global Uses of High Density Polyethylene}
\end{figure}
\begin{figure}[htbp]
\centering
\includegraphics[width=0.88\linewidth]{figures/fig_176.jpg}
\caption{EXHIBIT 10 - EXHIBIT 10: Polyethylene uses vary by grade (part 2: low density) Global Uses of Low Density Polyethylene}
\end{figure}
\begin{figure}[htbp]
\centering
\includegraphics[width=0.88\linewidth]{figures/fig_177.jpg}
\caption{EXHIBIT 11 - EXHIBIT 11: Polypropylene end uses are varied Polypropylene End Use by Application}
\end{figure}
\begin{figure}[htbp]
\centering
\includegraphics[width=0.88\linewidth]{figures/fig_178.jpg}
\caption{Exhibit 12 - EXHIBIT 12: Polypropylene production grows steadily over time ...}
\end{figure}
\begin{figure}[htbp]
\centering
\includegraphics[width=0.88\linewidth]{figures/fig_179.jpg}
\caption{EXHIBIT 13 - EXHIBIT 13: ... with no strong correlation with GDP growth rates.}
\end{figure}
\begin{figure}[htbp]
\centering
\includegraphics[width=0.88\linewidth]{figures/fig_180.jpg}
\caption{EXHIBIT 14 - EXHIBIT 14: The same applies for polyethylene (1/2) EXHIBIT 15: The same applies for polyethylene (2/2)}
\end{figure}
\begin{figure}[htbp]
\centering
\includegraphics[width=0.88\linewidth]{figures/fig_181.jpg}
\caption{EXHIBIT 14 - EXHIBIT 14: The same applies for polyethylene (1/2) EXHIBIT 15: The same applies for polyethylene (2/2)}
\end{figure}
\begin{figure}[htbp]
\centering
\includegraphics[width=0.88\linewidth]{figures/fig_182.jpg}
\caption{Exhibit 16 - EXHIBIT 16: Despite price volatility, polypropylene consumption continue to grow steadily over time...}
\end{figure}
\begin{figure}[htbp]
\centering
\includegraphics[width=0.88\linewidth]{figures/fig_183.jpg}
\caption{EXHIBIT 17 - EXHIBIT 17: ... so there's no strong correlation between consumption and price}
\end{figure}
\begin{figure}[htbp]
\centering
\includegraphics[width=0.88\linewidth]{figures/fig_184.jpg}
\caption{EXHIBIT 18 - EXHIBIT 18: The same applies to polyethylene (1/2)}
\end{figure}
\begin{figure}[htbp]
\centering
\includegraphics[width=0.88\linewidth]{figures/fig_185.jpg}
\caption{EXHIBIT 19 - EXHIBIT 19: The same applies to polyethylene (2/2)}
\end{figure}
\begin{figure}[htbp]
\centering
\includegraphics[width=0.88\linewidth]{figures/fig_186.jpg}
\caption{Exhibit 20 - EXHIBIT 20: Major Asian economies have dealt with little inflation in the last 20 years Inflation in major Asian economies vs Asian polyolefins consumption growth}
\end{figure}
\begin{figure}[htbp]
\centering
\includegraphics[width=0.88\linewidth]{figures/fig_187.jpg}
\caption{EXHIBIT 21 - EXHIBIT 21: Inflation therefore has little impact on polyolefin consumption growth rates Inflation in major Asian economies vs Asian polyolefins consumption growth correlation}
\end{figure}
\begin{figure}[htbp]
\centering
\includegraphics[width=0.88\linewidth]{figures/fig_188.jpg}
\caption{Exhibit 22 - EXHIBIT 22: Qatar LNG exports mainly go to Asia Qatar LNG exports by region}
\end{figure}
\begin{figure}[htbp]
\centering
\includegraphics[width=0.88\linewidth]{figures/fig_189.jpg}
\caption{EXHIBIT 23 - EXHIBIT 23: The US represents 65\% of European LNG imports European LNG imports by region}
\end{figure}
\begin{figure}[htbp]
\centering
\includegraphics[width=0.88\linewidth]{figures/fig_190.jpg}
\caption{EXHIBIT 24 - EXHIBIT 24: Most of the sites that declared force majeure are located in the Middle East and in Asia Sites that declared force majeure by country}
\end{figure}
\begin{figure}[htbp]
\centering
\includegraphics[width=0.88\linewidth]{figures/fig_191.jpg}
\caption{EXHIBIT 28 - EXHIBIT 28: Operating rates in Asia are understood to be falling}
\end{figure}
\begin{figure}[htbp]
\centering
\includegraphics[width=0.88\linewidth]{figures/fig_192.jpg}
\caption{Exhibit 24 - EXHIBIT 29: Asia is the leading global producer of Polypropylene... EXHIBIT 30: ... and polyethylene. EXHIBIT 31: We see slightly more than 8mt of PE capacity coming online in 2026, stemming from multiple regions}
\end{figure}
\begin{figure}[htbp]
\centering
\includegraphics[width=0.88\linewidth]{figures/fig_193.jpg}
\caption{EXHIBIT 29 - EXHIBIT 29: Asia is the leading global producer of Polypropylene... EXHIBIT 30: ... and polyethylene. EXHIBIT 31: We see slightly more than 8mt of PE capacity coming online in 2026, stemming from multiple regions}
\end{figure}
\begin{figure}[htbp]
\centering
\includegraphics[width=0.88\linewidth]{figures/fig_194.jpg}
\caption{EXHIBIT 30 - EXHIBIT 30: ... and polyethylene. EXHIBIT 31: We see slightly more than 8mt of PE capacity coming online in 2026, stemming from multiple regions}
\end{figure}
\begin{figure}[htbp]
\centering
\includegraphics[width=0.88\linewidth]{figures/fig_195.jpg}
\caption{EXHIBIT 32 - EXHIBIT 32: Slightly less than 6mt of polypropylene capacity is expected to come online in 2026 mainly from China, wider Asia and the Middle East}
\end{figure}
\begin{figure}[htbp]
\centering
\includegraphics[width=0.88\linewidth]{figures/fig_196.jpg}
\caption{Exhibit 33 - EXHIBIT 33: Our polyethylene model, back-tested, provides useful directional context for price forecasting (1/2)}
\end{figure}
\begin{figure}[htbp]
\centering
\includegraphics[width=0.88\linewidth]{figures/fig_197.jpg}
\caption{EXHIBIT 34 - EXHIBIT 34: Our polyethylene model, back-tested, provides useful directional context for price forecasting (2/2)}
\end{figure}
\begin{figure}[htbp]
\centering
\includegraphics[width=0.88\linewidth]{figures/fig_198.jpg}
\caption{EXHIBIT 35 - EXHIBIT 35: Our polypropylene model, back-tested, provides useful directional context for price forecasting (1/2)}
\end{figure}
\begin{figure}[htbp]
\centering
\includegraphics[width=0.88\linewidth]{figures/fig_199.jpg}
\caption{EXHIBIT 36 - EXHIBIT 36: Our polypropylene model, back-tested, provides useful directional context for price forecasting (2/2)}
\end{figure}
\begin{figure}[htbp]
\centering
\includegraphics[width=0.88\linewidth]{figures/fig_200.jpg}
\caption{Exhibit 37 - EXHIBIT 37: We expect the more muted demand growth rates of the early 2020s to continue}
\end{figure}
\begin{figure}[htbp]
\centering
\includegraphics[width=0.88\linewidth]{figures/fig_201.jpg}
\caption{EXHIBIT 38 - EXHIBIT 38: We now see less PE capacity coming online in 2026 as a result of capacity delays in the Middle East and in Asia... (1/2)}
\end{figure}
\begin{figure}[htbp]
\centering
\includegraphics[width=0.88\linewidth]{figures/fig_202.jpg}
\caption{EXHIBIT 39 - EXHIBIT 39: We now see less PE capacity coming online in 2026 as a result of capacity delays in the Middle East and in Asia... (2/2)}
\end{figure}
\begin{figure}[htbp]
\centering
\includegraphics[width=0.88\linewidth]{figures/fig_203.jpg}
\caption{EXHIBIT 40 - EXHIBIT 40: ... which we also see happening for Polypropylene (1/2)}
\end{figure}
\begin{figure}[htbp]
\centering
\includegraphics[width=0.88\linewidth]{figures/fig_204.jpg}
\caption{EXHIBIT 41 - EXHIBIT 41: ... which we also see happening for Polypropylene (2/2)}
\end{figure}
\begin{figure}[htbp]
\centering
\includegraphics[width=0.88\linewidth]{figures/fig_205.jpg}
\caption{EXHIBIT 42 - EXHIBIT 42: We increase our long-term operating rate assumptions for polyethylene}
\end{figure}
\begin{figure}[htbp]
\centering
\includegraphics[width=0.88\linewidth]{figures/fig_206.jpg}
\caption{EXHIBIT 43 - EXHIBIT 43: We increase our long-term operating rate assumptions for polypropylene}
\end{figure}
\begin{figure}[htbp]
\centering
\includegraphics[width=0.88\linewidth]{figures/fig_207.jpg}
\caption{Exhibit 44 - EXHIBIT 44: Our Global O\&G colleagues expect the oil price to settle at \textbackslash{}\$70/bbl}
\end{figure}
\begin{figure}[htbp]
\centering
\includegraphics[width=0.88\linewidth]{figures/fig_208.jpg}
\caption{Exhibit 45 - EXHIBIT 45: We expect GDP to slowly trend back to historical average growth rates}
\end{figure}
\begin{figure}[htbp]
\centering
\includegraphics[width=0.88\linewidth]{figures/fig_209.jpg}
\caption{Exhibit 46 - EXHIBIT 46: We increase our long-term PE prices from 5-10\% PE Avg. prices for each region are price weighted averages of HDPE, LLDPE and LDPE benchmarks EXHIBIT 47: We increase our long-term PP prices from 7-15\%}
\end{figure}
\begin{figure}[htbp]
\centering
\includegraphics[width=0.88\linewidth]{figures/fig_210.jpg}
\caption{EXHIBIT 46 - EXHIBIT 46: We increase our long-term PE prices from 5-10\% PE Avg. prices for each region are price weighted averages of HDPE, LLDPE and LDPE benchmarks EXHIBIT 47: We increase our long-term PP prices from 7-15\% \# BULL AND BEAR}
\end{figure}
\begin{figure}[htbp]
\centering
\includegraphics[width=0.88\linewidth]{figures/fig_211.jpg}
\caption{EXHIBIT 48 - EXHIBIT 48: Our polyolefin price scenarios}
\end{figure}
\begin{figure}[htbp]
\centering
\includegraphics[width=0.88\linewidth]{figures/fig_212.jpg}
\caption{EXHIBIT 49 - EXHIBIT 49: In a bull case scenario, we see polyolefin prices remain around current levels (1/2)}
\end{figure}
\begin{figure}[htbp]
\centering
\includegraphics[width=0.88\linewidth]{figures/fig_213.jpg}
\caption{EXHIBIT 50 - EXHIBIT 50: In a bull case scenario, we see polyolefin prices remain around current levels (2/2)}
\end{figure}
\begin{figure}[htbp]
\centering
\includegraphics[width=0.88\linewidth]{figures/fig_214.jpg}
\caption{Exhibit 53 - EXHIBIT 53: After the Middle-East conflict, the oil spread with Henry Hub spiked making gas-based polyolefin feedstocks more economic than oil-based ones Brent to Henry Hub spread - 2000/2026}
\end{figure}
\begin{figure}[htbp]
\centering
\includegraphics[width=0.88\linewidth]{figures/fig_215.jpg}
\caption{EXHIBIT 54 - EXHIBIT 54: The US Ethane advantage is high when oil prices are higher... Global Ethylene Variable Cost of Production by Feedstock on March 12 2025 with Brent Spot @ \textbackslash{}\$70.95 (\textbackslash{}\$/mt)}
\end{figure}
\begin{figure}[htbp]
\centering
\includegraphics[width=0.88\linewidth]{figures/fig_216.jpg}
\caption{EXHIBIT 55 - EXHIBIT 55: ... but that advantage declines when oil prices are lower Global Ethylene Variable Cost of Production by Feedstock on April 9 2025 with Brent Spot @ \textbackslash{}\$65.48 (\textbackslash{}\$/mt)}
\end{figure}
\begin{figure}[htbp]
\centering
\includegraphics[width=0.88\linewidth]{figures/fig_217.jpg}
\caption{EXHIBIT 56 - EXHIBIT 56: North America and the Middle East primarily rely on Ethane while Europe and China mainly produce petrochemicals with naphta Share of petrochemical feedstocks by region}
\end{figure}
\begin{figure}[htbp]
\centering
\includegraphics[width=0.88\linewidth]{figures/fig_218.jpg}
\caption{EXHIBIT 57 - EXHIBIT 58: On our house natural gas forecasts, the ethane advantage would disappear. But the on the market's... Henry Hub price forecast (\textbackslash{}\$/mmbtu)}
\end{figure}
\begin{figure}[htbp]
\centering
\includegraphics[width=0.88\linewidth]{figures/fig_219.jpg}
\caption{EXHIBIT 57 - EXHIBIT 58: On our house natural gas forecasts, the ethane advantage would disappear. But the on the market's... Henry Hub price forecast (\textbackslash{}\$/mmbtu)}
\end{figure}
\begin{figure}[htbp]
\centering
\includegraphics[width=0.88\linewidth]{figures/fig_220.jpg}
\caption{EXHIBIT 59 - EXHIBIT 59: We see being integrated, like BASF are, as an advantage in a world of higher petchem prices for longer \# We operate long and multiple-step value chains and sell products at every step in the value chain EO value chain as}
\end{figure}
{\small\color{refgray}\textbf{References:} [R009] 全球储能市场正在经历一场“去中国化”无法完成的供应链重构; [R014] 霍尔木兹海峡若重启，南非资产真正的赢家不是你想的那样; [R015] 市场真正低估的不是需求韧性，而是供给侧的永久性收缩}

\section{地产与消费}
\textbf{中国地产价格预期走弱但政策工具精准，新能源车爆款订单集中度上升重塑竞争格局。}

\begin{itemize}
  \item 上海以旧换新试点三个月成交523套，徐汇区占近九成，政策更像定向工具而非普惠激活器。
  \item 二手房中介信心指数连续三周下跌，看房量和认购量环比均降7\%，成交量稳定可能来自存量挂牌去化而非新增需求。
  \item 新能源车第21周订单环比增长16\%，小鹏GX发布12小时获近2.5万不可退订单，单周订单环比暴增433\%。
  \item 爆款车型的订单集中度急剧上升，竞争从谁卖得多转向谁能定义下一款爆款，对利润和供应链冲击更大。
  \item 新能源经销商平均折扣扩大到7.79\%，燃油车折扣率达19.67\%，价格战仍在持续。
\end{itemize}
\begin{figure}[htbp]
\centering
\includegraphics[width=0.88\linewidth]{figures/fig_148.jpg}
\caption{Exhibit 1 - Exhibit 1: Primary GFA sold last week was \$-4\textbackslash{}\%\$ wow and \$-5\textbackslash{}\%\$ yoy in c.75 cities}
\end{figure}
\begin{figure}[htbp]
\centering
\includegraphics[width=0.88\linewidth]{figures/fig_149.jpg}
\caption{Exhibit 2 - Exhibit 2: Primary GFA sold YTD on average was -14\% yoy in c.75 cities, and -14\%/-48\% vs. 2024/2023 level}
\end{figure}
\begin{figure}[htbp]
\centering
\includegraphics[width=0.88\linewidth]{figures/fig_150.jpg}
\caption{Exhibit 3 - Exhibit 3: Secondary GFA sold last week was +1\% wow and +11\% yoy in c.20 cities Average weekly volume of secondary property sales}
\end{figure}
\begin{figure}[htbp]
\centering
\includegraphics[width=0.88\linewidth]{figures/fig_151.jpg}
\caption{Exhibit 4 - Exhibit 4: Secondary GFA sold YTD was -1\% yoy in c.20 cities, while +23\%/+6\% vs. 2024/ secondary volume sold vs. 2022-25}
\end{figure}
\begin{figure}[htbp]
\centering
\includegraphics[width=0.88\linewidth]{figures/fig_152.jpg}
\caption{Exhibit 5 - Exhibit 5: Average CSI was -1.6pp wow and +2.3pp yoy Weekly Centraline Salesman Index (CSI) tracker in 5 cities}
\end{figure}
\begin{figure}[htbp]
\centering
\includegraphics[width=0.88\linewidth]{figures/fig_153.jpg}
\caption{Exhibit 6 - Exhibit 6: Average CAI was -0.4pp wow and -4.1pp yoy Weekly Centraline Seller Asking Index (CAI) tracker in 6 cities}
\end{figure}
\begin{figure}[htbp]
\centering
\includegraphics[width=0.88\linewidth]{figures/fig_154.jpg}
\caption{Exhibit 7 - Exhibit 7: Inventory balance was -0.2\% wow, -3.7\% from end-25 levels c.20 cities' total inventory breakdown by city tier (Indexed to Jan 2013)}
\end{figure}
\begin{figure}[htbp]
\centering
\includegraphics[width=0.88\linewidth]{figures/fig_155.jpg}
\caption{Exhibit 8 - Exhibit 8: Inventory month was -0.4\% wow, representing +0.2\% from end-25 levels c.20 cities' inventory months (12mth rolling) breakdown by city tier}
\end{figure}
\begin{figure}[htbp]
\centering
\includegraphics[width=0.88\linewidth]{figures/fig_156.jpg}
\caption{Exhibit 9 - Exhibit 9: MTD GSPC tracker implied monthly completions... GSPC tracker implied monthly GFA completion - based on GS float glass S-D model}
\end{figure}
\begin{figure}[htbp]
\centering
\includegraphics[width=0.88\linewidth]{figures/fig_157.jpg}
\caption{Exhibit 10 - Exhibit 10: ...suggesting completions at a high-teens \% yoy decline for MTD May-26 \% yoy change of GSPC - based on GS float glass S-D model}
\end{figure}
\begin{figure}[htbp]
\centering
\includegraphics[width=0.88\linewidth]{figures/fig_158.jpg}
\caption{Exhibit 2 - Exhibit 2: NEV dealer discount trend}
\end{figure}
\begin{figure}[htbp]
\centering
\includegraphics[width=0.88\linewidth]{figures/fig_159.jpg}
\caption{Exhibit 3 - Exhibit 3: ICE dealer discount trend}
\end{figure}
\begin{figure}[htbp]
\centering
\includegraphics[width=0.88\linewidth]{figures/fig_160.jpg}
\caption{Exhibit 3 - Exhibit 3: ICE dealer discount trend}
\end{figure}
\begin{figure}[htbp]
\centering
\includegraphics[width=0.88\linewidth]{figures/fig_161.jpg}
\caption{Exhibit 3 - Exhibit 3: ICE dealer discount trend}
\end{figure}
{\small\color{refgray}\textbf{References:} [R012] 市场真正低估的不是成交量，而是价格预期的自我强化; [R013] 小鹏GX单周2.5万不可退订单背后，市场真正低估的是“爆款能力”对竞争格局的重塑}

\section{区域市场}
\textbf{南非资产在霍尔木兹海峡重启情景下受益，但市场对供应链重构的二阶效应定价不足。}

\begin{itemize}
  \item 海峡重启后最值得关注的并非冲突期间跌幅最深的标的，而是贵金属、零售与金融板块。
  \item 油价即使重启后仍将粘性在100美元/桶以上，南非通胀路径和利率预期不会快速下行，仍需加息50-75个基点。
  \item 兰特走强将推动海外资金重新入场，国内主题重新启动，贵金属板块评分最高。
  \item 南非经济增长预期从1.2\%下调至1.0\%，重启交易的第一波驱动力不会是宏观基本面快速改善。
\end{itemize}
\begin{figure}[htbp]
\centering
\includegraphics[width=0.88\linewidth]{figures/fig_162.jpg}
\caption{Figure 2 - Figure 2: Gen. Retail may work for a short period but EPS compression is the risk}
\end{figure}
\begin{figure}[htbp]
\centering
\includegraphics[width=0.88\linewidth]{figures/fig_163.jpg}
\caption{Figure 4 - Figure 4: SA Domestic have de-rated vs EM broadening appeal for engagement to these names}
\end{figure}
\begin{figure}[htbp]
\centering
\includegraphics[width=0.88\linewidth]{figures/fig_164.jpg}
\caption{Figure 1 - Figure 1: ToT will be accommodated by an uplift from commodities which in turn turns positive for ZAR}
\end{figure}
\begin{figure}[htbp]
\centering
\includegraphics[width=0.88\linewidth]{figures/fig_165.jpg}
\caption{Figure 3 - Figure 3: This is in contrast to banks which are large, liquid and offer robust earnings}
\end{figure}
\begin{figure}[htbp]
\centering
\includegraphics[width=0.88\linewidth]{figures/fig_166.jpg}
\caption{Figure 5 - Figure 5: Discount to 10yr fwd. PE for JSE sectors}
\end{figure}
\begin{figure}[htbp]
\centering
\includegraphics[width=0.88\linewidth]{figures/fig_167.jpg}
\caption{Figure 6 - Figure 6: MSCI South Africa remains one of EM's highest beta markets}
\end{figure}
{\small\color{refgray}\textbf{References:} [R014] 霍尔木兹海峡若重启，南非资产真正的赢家不是你想的那样}

\section{风险偏好与监管}
\textbf{跨境监管重塑互联网券商估值逻辑，生命科学合作面临压力但长期韧性被低估。}

\begin{itemize}
  \item 中国证监会要求两年内清理存量不合规跨境业务，富途和老虎面临罚款与客户结构调整双重冲击。
  \item 罚款对富途2026年净利润影响约-3.93亿美元（-25\%），对老虎影响约-0.66亿美元（-60\%），客户结构调整是永久性的。
  \item 美国众议院建议将生物技术纳入对外投资禁止清单，但生命科学跨境合作的基本盘大概率不会脱钩。
  \item 中国生物科技公司从me-too向first-in-class原创分子授权输出的结构性转变，正在改变合作模式的定价权。
  \item 集采续约价格连续两次上行，冠脉支架中标价从730-848元升至839-949元，政策重心从压低价格转向稳定供应。
\end{itemize}
\begin{figure}[htbp]
\centering
\includegraphics[width=0.88\linewidth]{figures/fig_229.jpg}
\caption{Exhibit 2 - Exhibit 2: Earnings will primarily be impacted by regulatory fines but also the loss of revenues from remediating non-compliant mainland accounts. Exhibit 3: The cumulative impact represents -27\%/-38\% of FUTU's and TIGR's 2025 net}
\end{figure}
\begin{figure}[htbp]
\centering
\includegraphics[width=0.88\linewidth]{figures/fig_230.jpg}
\caption{Exhibit 5 - Exhibit 5: FUTU's competitive advantages include a significantly lower fee structure than traditional banks and brokers (3bps vs. \textbackslash{}\textasciitilde{}25bps/\textbackslash{}\textasciitilde{}7bps) and a willingness to invest in acquiring new clients Exhibit 6: FUTU has enjoyed rapi}
\end{figure}
\begin{figure}[htbp]
\centering
\includegraphics[width=0.88\linewidth]{figures/fig_231.jpg}
\caption{Exhibit 7 - Exhibit 7: Moreover, FUTU's market share has been continuously increasing in HK, from \$15\textbackslash{}\%\$ in 2021 to \$30\textbackslash{}\%\$ in 3Q25, and we expect it to reach \$47\textbackslash{}\%\$ by 2027.}
\end{figure}
\begin{figure}[htbp]
\centering
\includegraphics[width=0.88\linewidth]{figures/fig_232.jpg}
\caption{Exhibit 8 - \# 3. Higher Client Acquisition Cost (CACs) from International Expansion Both FUTU and TIGR have announced plans to expand into international markets. FUTU aims to enter South Korea as soon as possible, and while it has been cultivating the Japanese market for }
\end{figure}
\begin{figure}[htbp]
\centering
\includegraphics[width=0.88\linewidth]{figures/fig_233.jpg}
\caption{Exhibit 9 - Exhibit 9: ...and \$8\textbackslash{}\%\$ for TIGR}
\end{figure}
\begin{figure}[htbp]
\centering
\includegraphics[width=0.88\linewidth]{figures/fig_234.jpg}
\caption{Exhibit 10 - Exhibit 10: We have lowered AUM per client by -6\% for FUTU...}
\end{figure}
\begin{figure}[htbp]
\centering
\includegraphics[width=0.88\linewidth]{figures/fig_235.jpg}
\caption{Exhibit 11 - Exhibit 11: ...and by -3\% for TIGR}
\end{figure}
\begin{figure}[htbp]
\centering
\includegraphics[width=0.88\linewidth]{figures/fig_236.jpg}
\caption{Exhibit 12 - Neither FUTU nor TIGR have updated their new client growth guidance for 2026 (FUTU: 800k; TIGR: 150k). Even in a best case scenario that new client growth reaches guidance, higher corresponding CAC and lower AUM per new client will pressure revenues and profit}
\end{figure}
\begin{figure}[htbp]
\centering
\includegraphics[width=0.88\linewidth]{figures/fig_237.jpg}
\caption{Exhibit 14 - Exhibit 14: FUTU valuation range}
\end{figure}
\begin{figure}[htbp]
\centering
\includegraphics[width=0.88\linewidth]{figures/fig_238.jpg}
\caption{Exhibit 15 - Exhibit 15: FUTU 12M P/E vs. AUM growth}
\end{figure}
\begin{figure}[htbp]
\centering
\includegraphics[width=0.88\linewidth]{figures/fig_239.jpg}
\caption{Exhibit 16 - Exhibit 16: TIGR valuation range}
\end{figure}
\begin{figure}[htbp]
\centering
\includegraphics[width=0.88\linewidth]{figures/fig_240.jpg}
\caption{Exhibit 17 - Exhibit 17: FUTU 12M P/E vs. AUM growth}
\end{figure}
\begin{figure}[htbp]
\centering
\includegraphics[width=0.88\linewidth]{figures/fig_241.jpg}
\caption{Exhibit 1 - Exhibit 1: Price of each round of DES VBP continue to increase Price range of each round of DES VBP (in Rmb)}
\end{figure}
\begin{figure}[htbp]
\centering
\includegraphics[width=0.88\linewidth]{figures/fig_242.jpg}
\caption{Exhibit 2 - Exhibit 2: Submitted demand reached c.2.73mn units, implying a \$14\textbackslash{}\%\$ CAGR vs. the 2022 renewal First year guaranteed volume (mn units)}
\end{figure}
\begin{figure}[htbp]
\centering
\includegraphics[width=0.88\linewidth]{figures/fig_243.jpg}
\caption{Exhibit 3 - Exhibit 3: MNC submitted volume increased to c.41\% in the 2nd round renewal (vs. c.32\% in the first round) Total first-year procurement demand of major brands in the 1st and 2nd round renewal (units)}
\end{figure}
\begin{figure}[htbp]
\centering
\includegraphics[width=0.88\linewidth]{figures/fig_244.jpg}
\caption{Exhibit 6 - Exhibit 7: Accelerating progress in VBP since Oct 25; expecting additional announcements towards year end 2026}
\end{figure}
{\small\color{refgray}\textbf{References:} [R018] 监管的代价不是一次性罚款，而是重构客户结构的长期成本; [R019] 集采续约的定价信号：市场真正低估的不是需求，而是供给侧的再定价; [R020] 中美生命科学合作面临新压力，但真正被低估的是行业的韧性再定价}

\section{信用与资产上链}
\textbf{EM信用债曲线趋平反映期限溢价系统性重定价，代币化RWA逆势增长驱动金融基础设施变革。}

\begin{itemize}
  \item EM Aggregate 10s30s利差曲线当前68个基点，较一年前收窄7个基点，曲线中枢下移反映全球资本对新兴市场长期信用风险评估框架调整。
  \item EM投资级曲线斜率61个基点，高收益达83个基点，埃及以159个基点成为最陡峭曲线，印尼以20个基点最平坦。
  \item 代币化RWA总市值突破510亿美元，年内增长42\%，同期加密市场下跌约15\%，机构资金正在寻找更稳健的链上资产载体。
  \item 私人信贷占RWA总市值44\%，美国国债占30\%，以太坊和Provenance链承载约75\%的代币化资产。
  \item Hyperliquid上RWA衍生品已占平台永续合约交易量35\%，股票代币化机构跑步入场，DTCC联合50多家金融机构推进。
\end{itemize}
\begin{figure}[htbp]
\centering
\includegraphics[width=0.88\linewidth]{figures/fig_001.jpg}
\caption{Figure 1 - Figure 1: EM USD Aggregate 10s30s}
\end{figure}
\begin{figure}[htbp]
\centering
\includegraphics[width=0.88\linewidth]{figures/fig_002.jpg}
\caption{Figure 2 - Figure 2: UST 10s30s yield curve slope}
\end{figure}
\begin{figure}[htbp]
\centering
\includegraphics[width=0.88\linewidth]{figures/fig_003.jpg}
\caption{Figure 3 - Figure 3: Steepest and flattest 10s30s curves}
\end{figure}
\begin{figure}[htbp]
\centering
\includegraphics[width=0.88\linewidth]{figures/fig_004.jpg}
\caption{Figure 4 - Figure 4: Largest 1w changes in 10s30s}
\end{figure}
\begin{figure}[htbp]
\centering
\includegraphics[width=0.88\linewidth]{figures/fig_005.jpg}
\caption{Figure 5 - Figure 6: 10s30s curves vs. the 10y level relationship}
\end{figure}
\begin{figure}[htbp]
\centering
\includegraphics[width=0.88\linewidth]{figures/fig_006.jpg}
\caption{Figure 7 - Figure 7: 1w change in aggregate 10s30s vs. 1w change in aggregate 10 spreads}
\end{figure}
\begin{figure}[htbp]
\centering
\includegraphics[width=0.88\linewidth]{figures/fig_007.jpg}
\caption{Figure 8 - Figure 8: 1w change in 10s30s vs. 1w change in 10y spreads by issuer Figure 9: Historical EM aggregate 10s30s spread curve slope}
\end{figure}
\begin{figure}[htbp]
\centering
\includegraphics[width=0.88\linewidth]{figures/fig_008.jpg}
\caption{Figure 8 - Figure 8: 1w change in 10s30s vs. 1w change in 10y spreads by issuer Figure 9: Historical EM aggregate 10s30s spread curve slope}
\end{figure}
\begin{figure}[htbp]
\centering
\includegraphics[width=0.88\linewidth]{figures/fig_009.jpg}
\caption{Figure 11 - Figure 11: Historical EMBIG vs. CEMBI spread curve slope}
\end{figure}
\begin{figure}[htbp]
\centering
\includegraphics[width=0.88\linewidth]{figures/fig_010.jpg}
\caption{Figure 13 - Figure 13: Historical sovereign vs. quasi-sovereign curve slope}
\end{figure}
\begin{figure}[htbp]
\centering
\includegraphics[width=0.88\linewidth]{figures/fig_011.jpg}
\caption{Figure 15 - Figure 15: Historical EMBIGD curve slope by region}
\end{figure}
\begin{figure}[htbp]
\centering
\includegraphics[width=0.88\linewidth]{figures/fig_012.jpg}
\caption{Figure 10 - Figure 10: Historical US Treasury 10s30s yield curve slope}
\end{figure}
\begin{figure}[htbp]
\centering
\includegraphics[width=0.88\linewidth]{figures/fig_013.jpg}
\caption{Figure 12 - Figure 12: Historical EM IG vs. HY spread curve slope}
\end{figure}
\begin{figure}[htbp]
\centering
\includegraphics[width=0.88\linewidth]{figures/fig_014.jpg}
\caption{Figure 14 - Figure 14: EM corporate vs. US HG corporate spread curve slope}
\end{figure}
\begin{figure}[htbp]
\centering
\includegraphics[width=0.88\linewidth]{figures/fig_015.jpg}
\caption{Figure 16 - Figure 16: Historical CEMBI curve slope by region}
\end{figure}
\begin{figure}[htbp]
\centering
\includegraphics[width=0.88\linewidth]{figures/fig_016.jpg}
\caption{Figure 17 - Figure 17: 10s30s spread curve slopes vs. 10y spread}
\end{figure}
\begin{figure}[htbp]
\centering
\includegraphics[width=0.88\linewidth]{figures/fig_017.jpg}
\caption{Figure 18 - Figure 18: Asia issuers 10s30s spread curve}
\end{figure}
\begin{figure}[htbp]
\centering
\includegraphics[width=0.88\linewidth]{figures/fig_018.jpg}
\caption{Figure 19 - Figure 19: CEEMEA issuers}
\end{figure}
\begin{figure}[htbp]
\centering
\includegraphics[width=0.88\linewidth]{figures/fig_019.jpg}
\caption{Figure 20 - Figure 20: Latin America issuers}
\end{figure}
\begin{figure}[htbp]
\centering
\includegraphics[width=0.88\linewidth]{figures/fig_020.jpg}
\caption{Figure 21 - Figure 21: 10s30s spread curve slopes vs. 10y spread}
\end{figure}
\begin{figure}[htbp]
\centering
\includegraphics[width=0.88\linewidth]{figures/fig_021.jpg}
\caption{Figure 22 - Figure 22: EMBIG sovereign issuers}
\end{figure}
\begin{figure}[htbp]
\centering
\includegraphics[width=0.88\linewidth]{figures/fig_022.jpg}
\caption{Figure 23 - Figure 23: EMBIG quasi-sovereign issuers Figure 24: CEMBI issuers \# Slope versus Credit Rating Relationship Figure 25: 10s30s spread curve slopes vs. average credit rating}
\end{figure}
\begin{figure}[htbp]
\centering
\includegraphics[width=0.88\linewidth]{figures/fig_023.jpg}
\caption{Figure 23 - Figure 23: EMBIG quasi-sovereign issuers Figure 24: CEMBI issuers \# Slope versus Credit Rating Relationship Figure 25: 10s30s spread curve slopes vs. average credit rating}
\end{figure}
\begin{figure}[htbp]
\centering
\includegraphics[width=0.88\linewidth]{figures/fig_024.jpg}
\caption{Figure 24 - Figure 24: CEMBI issuers \# Slope versus Credit Rating Relationship Figure 25: 10s30s spread curve slopes vs. average credit rating}
\end{figure}
\begin{figure}[htbp]
\centering
\includegraphics[width=0.88\linewidth]{figures/fig_025.jpg}
\caption{Figure 26 - Figure 26: Asia issuers 10s30s spread curve}
\end{figure}
\begin{figure}[htbp]
\centering
\includegraphics[width=0.88\linewidth]{figures/fig_026.jpg}
\caption{Figure 27 - Figure 27: CEEMEA issuers Figure 28: Latin America issuers}
\end{figure}
\begin{figure}[htbp]
\centering
\includegraphics[width=0.88\linewidth]{figures/fig_027.jpg}
\caption{Figure 27 - Figure 27: CEEMEA issuers Figure 28: Latin America issuers}
\end{figure}
\begin{figure}[htbp]
\centering
\includegraphics[width=0.88\linewidth]{figures/fig_028.jpg}
\caption{Figure 29 - Figure 29: Steepest 10s30s spread curves across EM sovereign and corporate credit Figure 30: Steepest 10s30s curves vs. the 10y spread}
\end{figure}
\begin{figure}[htbp]
\centering
\includegraphics[width=0.88\linewidth]{figures/fig_029.jpg}
\caption{Figure 31 - Figure 31: Steepest 10s30s curves and 1w change}
\end{figure}
\begin{figure}[htbp]
\centering
\includegraphics[width=0.88\linewidth]{figures/fig_030.jpg}
\caption{Figure 32 - Figure 32: Flattest 10s30s spread curves across EM sovereign and corporate credit Figure 33: Flattest 10s30s curves vs. the 10y spread}
\end{figure}
\begin{figure}[htbp]
\centering
\includegraphics[width=0.88\linewidth]{figures/fig_031.jpg}
\caption{Figure 34 - Figure 34: Flattest 10s30s curves and 1w change}
\end{figure}
\begin{figure}[htbp]
\centering
\includegraphics[width=0.88\linewidth]{figures/fig_032.jpg}
\caption{Figure 35 - Figure 35: Largest 1w steepening across 10s30s spread curves Figure 36: 1w change in 10s30s curves vs. the current 10s30s}
\end{figure}
\begin{figure}[htbp]
\centering
\includegraphics[width=0.88\linewidth]{figures/fig_033.jpg}
\caption{Figure 37 - Figure 37: Largest 1w steepening and current 10s30s slope}
\end{figure}
\begin{figure}[htbp]
\centering
\includegraphics[width=0.88\linewidth]{figures/fig_034.jpg}
\caption{Figure 38 - Figure 38: Largest 1w flattening across 10s30s spread curves Figure 39: 1w change in 10s30s curves vs. the current 10s30s}
\end{figure}
\begin{figure}[htbp]
\centering
\includegraphics[width=0.88\linewidth]{figures/fig_035.jpg}
\caption{Figure 40 - Figure 40: Largest 1w flattening and current 10s30s slope}
\end{figure}
\begin{figure}[htbp]
\centering
\includegraphics[width=0.88\linewidth]{figures/fig_068.jpg}
\caption{Exhibit 8 - EXHIBIT 1: Tokenized RWA - Market Cap (\textbackslash{}\$Bn)}
\end{figure}
\begin{figure}[htbp]
\centering
\includegraphics[width=0.88\linewidth]{figures/fig_069.jpg}
\caption{EXHIBIT 2 - EXHIBIT 2: Tokenized RWA market cap split by asset class Tokenized RWA - Asset Class Split (\%)}
\end{figure}
\begin{figure}[htbp]
\centering
\includegraphics[width=0.88\linewidth]{figures/fig_070.jpg}
\caption{EXHIBIT 3 - EXHIBIT 3: Commodities and Indices dominate HIP-3 trading volumes HIP-3 Trading Volumes - By Assets (\textbackslash{}\$Bn)}
\end{figure}
\begin{figure}[htbp]
\centering
\includegraphics[width=0.88\linewidth]{figures/fig_071.jpg}
\caption{EXHIBIT 4 - EXHIBIT 4: HIP-3 represented \$35\textbackslash{}\%\$ of hyperliquid perp volumes in April'26 HIP-3 Volume Market Share of Hyperliquid Perpetuals (\%) - April'26}
\end{figure}
\begin{figure}[htbp]
\centering
\includegraphics[width=0.88\linewidth]{figures/fig_072.jpg}
\caption{EXHIBIT 5 - EXHIBIT 5: Tokenized RWA split by blockchain network Tokenized RWA - Blockchain network split (\%)}
\end{figure}
\begin{figure}[htbp]
\centering
\includegraphics[width=0.88\linewidth]{figures/fig_073.jpg}
\caption{EXHIBIT 6 - EXHIBIT 6: Top RWA tokenization platform EXHIBIT 7: Tokenized real world asset holders RWA Asset Holders (in '000)}
\end{figure}
\begin{figure}[htbp]
\centering
\includegraphics[width=0.88\linewidth]{figures/fig_074.jpg}
\caption{EXHIBIT 8 - EXHIBIT 8: FIGR - Consumer loan volumes Consumer loan volume (\$Bn)}
\end{figure}
\begin{figure}[htbp]
\centering
\includegraphics[width=0.88\linewidth]{figures/fig_075.jpg}
\caption{EXHIBIT 9 - EXHIBIT 9: Tokenized Equities - Monthly Transfer Volumes Tokenized Equity - Transfer Volumes (\$Bn)}
\end{figure}
{\small\color{refgray}\textbf{References:} [R003] EM信用债曲线正在传递一个被低估的信号; [R008] 市场低估的不是加密货币，而是资产上链对金融基础设施的重构}

\section*{结语}
综合来看，市场当前定价的核心偏差在于：对政策方向逆转（美联储加息、中国集采续约涨价）、结构性供给收缩（化工、储能、AI电力链路）以及统计口径扭曲（AI对GDP拉动）的认知不足。代币化与人民币国际化等制度性变革正在重塑金融基础设施，但短期地缘政治与监管风险仍需警惕。
\vfill
{\small\color{refgray}Personal reading notes and learning share only. Not investment advice.}
\end{document}
